\section{Why the Composition Fails}
\label{sec:mechanism}

Four accounts, labelled by evidential status: what is measured, and what is
hypothesis consistent with the measurement.

\subsection{Routing divides the data, not the update budget}
\label{sec:dilution}

\textbf{[Ordering MEASURED; arithmetic MEASURED; mechanism HYPOTHESIS.]} The
obvious account of the ordering is capacity dilution: MoR's single block sees
100\% of tokens while MoE and MoRE split the same FFN budget across six experts,
so each expert is trained less. \textbf{That account is arithmetically false
here, and ruling it out is what makes the surviving hypothesis worth stating.}

The budgets are matched by construction, and the matching cancels. Each MoE/MoRE
expert FFN holds $2 \cdot 4 \cdot d^2 = 524{,}288$ weights and MoR's single FFN
holds $2 \cdot 24 \cdot d^2 = 3{,}145{,}728$, exactly six times as many
(\S\ref{sec:engine}). A run sees $404.2$M tokens. So:

\begin{center}
\begin{tabular}{lccc}
\toprule
per FFN parameter & MoE expert & MoR block & MoRE expert \\
\midrule
tokens seen        & 128.5 & 128.5 & 128.5 \\
FFN applications   & 128   & 865   & \textbf{894} \\
\bottomrule
\end{tabular}
\end{center}

Tokens per parameter is \emph{identical} across all three arms---the token split
and the parameter split are both by $E$, so they cancel exactly. And once
recursion is counted, applications per parameter slightly \emph{favours} MoRE
(894) over MoR (865), because MoRE runs marginally deeper. Whatever costs MoRE
0.12 nats against MoR, it is not that its parameters are trained less.

That comparison also explains the easy half of the ordering. MoRE gets
$894/128 \approx 7\times$ the FFN applications per parameter that MoE gets, which
is simply what recursion buys, and MoRE beats MoE by 0.311 nats. The recursive
multiply is real and it is the entire source of MoRE's advantage over MoE.

\textbf{What remains, once update count is excluded, is the \emph{diversity} of
what each parameter block sees.} MoR's block is exposed to the whole token
distribution; each MoE/MoRE expert is exposed only to the router-selected slice
that reaches it, which is both narrower and systematically biased, since the
router selects it. Structure shared across the whole distribution---the bulk of
what a language model must learn at this scale---can be learned once by MoR's
block but has to be re-learned independently inside all six experts. Six partial
copies of shared structure fit in the same parameter budget as one complete copy
only by making each copy worse.

This is a hypothesis, not an isolated effect, and we label it as one: separating
it from the remaining differences needs a width sweep or a controlled
manipulation of routing diversity, and we ran neither. It earns its place because
it is the only account left after the arithmetic above excludes the update-count
version, because it predicts the stratified result---MoRE's advantage over MoE
grows with token difficulty (\S\ref{sec:stratified}), which is where an
incompletely-learned shared structure should hurt most---and because it generates
the falsifiable scale prediction in \S\ref{sec:limitations}: as the corpus grows,
each expert's slice eventually contains enough of the shared structure for the
re-learning penalty to amortise, so the MoR$-$MoRE gap should close with data
scale at fixed parameters.

\subsection{Depth saturation is a data-scale effect, not an architectural limit}
\label{sec:saturation}

\textbf{[MEASURED, four points at fixed hyperparameters.]} MoRE's halting sits at
the ceiling on the full corpus (97.3\% forced-exit, mean depth 6.955 of 7). It is
tempting to read that as the architecture being unable to learn adaptive depth. A
subset sweep falsifies that reading.

\begin{center}
\begin{tabular}{lcc}
\toprule
training data & mean depth & forced-exit rate \\
\midrule
2.5\% of the corpus & 2.43 & $\approx 0$ \\
10\% & 1.92 & $\approx 0$ \\
40\% & 1.97 & $\approx 0$ \\
100\% (canonical) & 6.96 & 0.973 \\
\bottomrule
\end{tabular}
\end{center}

At every reduced-data point the same model, at the same ponder weight, allocates
depth adaptively---exit rates near zero, mean depth between 1.9 and 2.4, inside a
budget of 7. The collapse appears only at full scale. \textbf{So the failure is a
calibration finding, not an architectural one}: the ponder penalty that produces
adaptive depth on a small corpus cannot hold depth down at this one.

We state the limit of that claim precisely, because the sweep is exploratory
(single seed, subset fraction below 1, and it may never be quoted as a canonical
result). The subset axis is \emph{not} cleanly monotone---2.5\% gives 2.43 but 40\%
gives 1.97---so the honest phrasing is that \emph{saturation emerges somewhere
between 40\% and 100\% of the data}, not that depth is monotone in training steps.
The mechanism is consistent: the ponder term is linear in expected depth, while the
task gradient grows as the model fits, so the fixed penalty is outpaced. Once the
halting sigmoid saturates, its gradient vanishes and the linear-penalty ACT scheme
cannot recover on its own \citep{graves2016act}. \citet{banino2021pondernet}
replaces the linear penalty with a KL divergence against a geometric prior, which
keeps a gradient on the halt head near saturation, and is the principled fix; we
did not run it, and it is named as future work rather than as a result.

\subsection{The two mechanisms compete for one budget}
\label{sec:competition}

\textbf{[MEASURED on a calibration grid; the canonical-corpus form is measured on
the balance axis.]} This is the paper's central claim, and it is the one that
explains why the composition does not merely fail to add but subtracts.

Depth and specialisation are not independent knobs in MoRE the way they are in MoR
and MoE separately. On the calibration grid, recovering adaptive depth by raising
the halting weight destroys expert differentiation:

\begin{center}
\begin{tabular}{lccc}
\toprule
halting weight & mean depth & AMI & val loss \\
\midrule
0.001 (frozen) & 2.698 & \textbf{0.2851} & 5.2352 \\
0.107 & 1.476 & \textbf{0.0025} & 5.3346 \\
0.500 & 1.072 & 0.0207 & 5.3306 \\
\bottomrule
\end{tabular}
\end{center}

AMI of 0.0025 is statistically indistinguishable from the shuffled-label control.
The mechanism is direct: fewer recursion steps means fewer expert applications per
token, which means fewer opportunities for experts to differentiate. \textbf{Tune
toward adaptive depth and specialisation collapses; preserve specialisation and
depth saturates.} Both mechanisms draw on the same per-token compute budget, and at
this parameter and data budget neither gets enough. (This grid is on a smaller dev
corpus at one seed; it is a mechanism demonstration and carries that caveat.)

The competition has a second face on the canonical corpus, along a different axis.
The routing balance term is load-bearing and cannot be traded away for depth
either: setting it to zero \emph{collapses} the expert load onto a single expert
(load entropy 0.0001) and drives AMI to zero. A near-uniform per-token router
distribution coexisting with balanced load and above-chance specialisation is
therefore not a pathology to be tuned out---and we note that we initially read it
as one before the sweep refuted us (\S\ref{sec:honesty}).

Specialisation itself is real, and measured against per-run controls rather than
against chance. MoRE's AMI is $0.1803 \pm 0.0246$ against a permutation control of
$0.0249$, a controlled z of $13.43$; Hungarian-matched accuracy is $0.3592 \pm
0.0271$ against a control of $0.2464$. MoE's AMI is $0.1637 \pm 0.0478$ against a
control of $0.0364$ ($z = 6.82$), Hungarian $0.3616$ against $0.2966$. Both
six-expert arms find substantial above-chance structure---and, read together, the
two arms are not distinguishable from each other. Routing quality is approximately
the same with recursion as without it, which is the routing half of the
interference result: recursion does not improve routing, and routing does not
preserve adaptive depth, in the same architecture.

\subsection{A plausible mechanism, refuted}
\label{sec:trapping}

\textbf{[MEASURED; hypothesis falsified.]} The natural explanation for MoRE's
collapse is \emph{subspace trapping}: routing a token once confines its recursion
to a single expert's parameter subspace, and repeated application of one block
within that subspace drives the state toward a fixed point rather than refining it.
It is an appealing account and it is wrong. We tested it three ways.

\begin{center}
\begin{tabular}{lccc}
\toprule
probe & MoE & MoR & MoRE \\
\midrule
final-state participation ratio & $41.5 \pm 1.6$ & $\mathbf{70.3 \pm 5.6}$ & $57.3 \pm 2.8$ \\
between-expert mean cosine & $0.630$ & \textsc{n/a} & $0.916$ \\
\bottomrule
\end{tabular}
\end{center}

\textbf{(i) Trapping predicts lower participation ratio in the composed arm;
it is higher.} If recursion were collapsing the state onto a low-dimensional
subspace, the composed model's final-state PR should fall below MoE's, since MoE
has no recursion at all. It does not: MoRE's PR is $57.3$ against MoE's $41.5$,
and MoR's---with no routing---is higher still at $70.3$. Recursion \emph{expands}
the effective dimensionality, and routing contracts it somewhat without collapsing
it. The ordering refutes the prediction.

\textbf{(ii) Trapping predicts oscillating or fixed-point dynamics; velocity
decays smoothly.} Measuring the L2 norm of consecutive post-norm states across
depth, both recursive arms decay monotonically---MoR from 11.06 at step 2 to 5.71
at step 7, MoRE from 9.09 to 3.49---which is the signature of settling, not of
vibrating around an attractor.

\textbf{(iii) Trapping predicts that successive updates rescale the state;
they are near-orthogonal.} The mean cosine between a step's update and the state
it updates from is negative in both arms ($-0.26$ to $-0.06$), i.e. updates push
into new directions rather than amplifying existing ones. A rescaling collapse
would show positive collinearity approaching $+1$.

The recursion in this architecture is dynamically healthy: it travels, it settles,
and it adds dimensions. The failure is not a broken iterated map. That rules out
the most attractive explanation and leaves the budget competition
(\S\ref{sec:competition}) as the account the data supports---which is a stronger
position than an unexplained null, because we can say which of the obvious
mechanisms it is \emph{not}.

